%\documentclass[twocolumn]{article}\
\documentclass{article}
\usepackage{amsmath, amssymb, cancel, mathtools, bm}
\usepackage[left=.5in, right=.5in, top=1in, bottom=1in]{geometry}
\usepackage[most]{tcolorbox}
\usepackage[skip=1em,indent=0pt]{parskip}
\setlength{\parindent}{0pt}
\newcommand{\svec}[1]{\underset{\sim}{\bm{#1}}} % Vector symbol (tilde under vector)
\DeclareMathOperator{\EX}{\mathbb{E}} % Expected Value symbol
\DeclareMathOperator{\ext}{\EX_{\theta}} % Expected Value symbol
\DeclareMathOperator{\vart}{Var_{\theta}} % Expected Value symbol
\newcommand{\indep}{\perp\!\!\!\!\perp} % Independence symbol
\newcommand{\real}{\mathbb{R}} % Simplified 'Reals' indicator
\newcommand{\pto}{\overset{P}{\to}}
\newcommand{\asto}{\overset{a.s.}{\to}}
\newcommand{\rthto}{\overset{L^r}{\to}}
\newcommand{\dto}{\overset{D}{\to}}
\newcommand{\xtb}{x^{\top}\svec{\beta}}
\newcommand{\xitb}{x_i^{\top}\svec{\beta}}
% Force all aggregate symbols to always put values above/below
\let\Oldint=\int
\let\Oldsum=\sum
\let\Oldprod=\prod
\let\Oldbigcup=\bigcup
\let\Oldbigcap=\bigcap
\let\Oldlim=\lim
\renewcommand{\int}{\Oldint\limits} 
\renewcommand{\sum}{\Oldsum\limits} 
\renewcommand{\prod}{\Oldprod\limits} 
\renewcommand{\bigcup}{\Oldbigcup\limits} 
\renewcommand{\bigcap}{\Oldbigcap\limits} 
\renewcommand{\lim}{\Oldlim\limits} 
\newcommand{\infint}{\int_{-\infty}^{\infty}}
\newcommand{\iiddist}{\overset{\mathrm{iid}}{\sim}}
\newcommand{\norm}[1]{\left\lVert#1\right\rVert}
\newcommand{\boldred}[1]{\textbf{\textcolor{red}{#1}}}


\begin{document}

\section{MLE of $\sigma^2$ from Multi-variate Normal}

Likelihood of the Multi-variate normal: 

$$L(\svec{\beta}, \sigma^2 | \svec{y}) = (2\pi)^{-n/2}(\sigma^2)^{-n/2} e^{-\frac{1}{2\sigma^2}(\svec{y} - x\svec{\beta})^{\top}(\svec{y} - x\svec{\beta})}$$

We can then write the log-likelihood:

$$l(\svec{\beta}, \sigma^2 | \svec{y}) = -\frac{n}{2}ln(2\pi) - \frac{n}{2}ln(\sigma^2) - \frac{1}{2\sigma^2}(\svec{y} - x\svec{\beta})^{\top}(\svec{y} - x\svec{\beta})$$

We then take the partial derivative with respect to $\sigma^2$ and set equal to $0$:

$$\frac{\partial l}{\partial \sigma^2} = 0 - \frac{n}{2\hat{\sigma}^2} + \frac{1}{2(\hat{\sigma}^2)^2}(\svec{y} - x\svec{\beta})^{\top}(\svec{y} - x\svec{\beta}) = 0$$

$$\implies n\hat{\sigma}^2 = (\svec{y} - x\svec{\beta})^{\top}(\svec{y} - x\svec{\beta}))$$

$$\implies \hat{\sigma}^2 = \frac{(\svec{y} - x\svec{\beta})^{\top}(\svec{y} - x\svec{\beta})}{n}$$

The $(\svec{y} - x\svec{\beta})^{\top}(\svec{y} - x\svec{\beta})$ is actually the same as the SSE, so this can be rewritten as:

$$\hat{\sigma}^2 = \frac{SSE}{n}$$


\section{MLE of $\hat{\svec{\beta}}$ from Multi-variate Normal}

$$L = (2\pi)^{-n/2}(\sigma^2)^{-n/2}e^{(-\frac{1}{2\sigma^2}(y-x\beta)^{\top}(y-x\beta))}$$

$$l = -\frac{n}{2}ln(2\pi) - \frac{1}{2}ln(|\Sigma|) - \frac{1}{2}(y-x\beta)^{\top}\Sigma^{-1}(y-x\beta)$$

$$
\begin{aligned}
(y-x\beta)^{\top}\Sigma^{-1}(y-x\beta) 
& = y^{\top}\Sigma^{-1}y - y^{\top}\Sigma^{-1}x\beta - (x\beta)^{\top}\Sigma^{-1}y + (x\beta)^{\top}\Sigma^{-1}x\beta \\
& = y^{\top}\Sigma^{-1}y - 2\beta^{\top}x^{\top}\Sigma^{-1}y + \beta^{\top}x^{\top}\Sigma^{-1}(x\beta)
\end{aligned}$$

$$\frac{\partial}{\partial \beta} -\frac{1}{2}(y-x\beta)^{\top}\Sigma^{-1}(y-x\beta) = -\frac{1}{2}[0 - 2x^{\top}\Sigma^{-1}y + 2x^{\top}\Sigma^{-1}x\beta] = x^{\top}\Sigma^{-1}y - x^{\top}\Sigma^{-1}x\beta$$

Set it equal to 0: 

$$
\begin{aligned}
0 &= x^{\top}\Sigma^{-1}y - x^{\top}\Sigma^{-1}x\hat{\beta} \\
x^{\top}\Sigma^{-1}x\hat{\beta} &= x^{\top}\Sigma^{-1}y \\
\hat{\beta} &= (x^{\top}\Sigma^{-1}x)^{-1}x^{\top}\Sigma^{-1}y
\end{aligned}$$

\end{document}